\documentclass[11pt,a4paper]{article}
\usepackage[utf8]{inputenc}
\usepackage[T1]{fontenc}
\usepackage[catalan]{babel}
\usepackage{lmodern}
\usepackage[sfdefault,lining]{FiraSans}
\usepackage[a4paper,top=2.3cm,bottom=2.2cm,left=1.8cm,right=1.8cm,headheight=1.3cm,headsep=0.4cm,footskip=1cm]{geometry}
\usepackage{amsmath,amssymb}
\allowdisplaybreaks
\usepackage{icomma}
\usepackage{graphicx}
\usepackage[table]{xcolor}
\usepackage{tikz}
\usepackage{fontawesome5}
\usepackage{ccicons}
\usepackage{fancyhdr}
\usepackage{lastpage}
\usepackage{tasks}
\usepackage{needspace}
\usepackage[most]{tcolorbox}
\usepackage{array}
\usepackage[hidelinks,pdftitle={Fitxa d'exercicis de la unitat 5: Límits i continuïtat. 2n de Batxillerat MII},pdfsubject={Matemàtiques II},pdfkeywords={límits, indeterminacions, infinitèsims equivalents, límits trigonomètrics, paràmetres, continuïtat, funcions a trossos, teorema de Bolzano, asímptotes, PAU}]{hyperref}

\definecolor{navy}{HTML}{1B2748}
\definecolor{accent}{HTML}{E8603A}
\definecolor{teal}{HTML}{1F9E90}
\definecolor{soft}{HTML}{EEF1F8}
\definecolor{muted}{HTML}{5D6685}

\setlength{\parindent}{0pt}
\setlength{\parskip}{3pt}
\settasks{label=\alph*), label-format=\bfseries\color{navy}, label-width=1.3em, item-indent=1.8em, column-sep=0.6em, after-item-skip=5pt, before-skip=3pt, after-skip=4pt}

% ---------- capçalera i peu ----------
\newcommand{\llicencia}{\href{https://creativecommons.org/licenses/by-nc-sa/4.0/deed.ca}{\ccbyncsa}\ \ Llicència Creative Commons Reconeixement-NoComercial-CompartirIgual 4.0 Internacional (CC BY-NC-SA 4.0)}
\pagestyle{fancy}
\fancyhf{}
\fancyhead[L]{\raisebox{-0.35\height}{\includegraphics[height=0.95cm]{logo-gva.png}}}
\fancyhead[C]{\small\color{muted}Fitxa d'exercicis · Unitat 5: Límits i continuïtat · 2n Batxillerat MII}
\fancyhead[R]{\raisebox{-0.35\height}{\includegraphics[height=0.95cm]{logo-vermellar.png}}}
\fancyfoot[L]{\scriptsize\color{muted}\llicencia}
\fancyfoot[R]{\small\color{muted}\thepage}
\renewcommand{\headrulewidth}{0.4pt}
\renewcommand{\headrule}{\color{muted!40}\hrule height \headrulewidth}
\fancypagestyle{primera}{\fancyhf{}\renewcommand{\headrulewidth}{0pt}
  \fancyfoot[L]{\scriptsize\color{muted}\llicencia}
  \fancyfoot[R]{\small\color{muted}\thepage}}

% ---------- nivells i exercicis ----------
\newcommand{\estrelles}[1]{{\color{accent}\ifcase#1\or\faStar\faStar[regular]\faStar[regular]\or\faStar\faStar\faStar[regular]\or\faStar\faStar\faStar\fi}}
\newcounter{exer}
\newcommand{\ex}[1]{\par\needspace{4\baselineskip}\medskip\refstepcounter{exer}\noindent{\color{navy}\bfseries\theexer.}\ \ }

\newcommand{\seccio}[2]{\par\needspace{10\baselineskip}\bigskip
  \begin{tcolorbox}[colback=navy,colframe=navy,coltext=white,arc=2mm,boxrule=0pt,left=3mm,right=3mm,top=1.2mm,bottom=1.2mm]
  {\large\bfseries #1\quad #2}
  \end{tcolorbox}\vspace{-2pt}}


\setlength{\emergencystretch}{3em}

% quadre de recordatori (només a la fitxa d'exercicis)
\newcommand{\recorda}[1]{\par\smallskip\begin{tcolorbox}[colback=accent!7,colframe=accent!55,arc=2mm,boxrule=0.5pt,left=2.5mm,right=2.5mm,top=1mm,bottom=1mm]\small{\color{accent}\faLightbulb[regular]}\ \textbf{Recorda:}\ #1\end{tcolorbox}}
% etiqueta dels exercicis de la PAU d'altres comunitats
\newcommand{\pau}[1]{{\color{teal}\small\bfseries\faUniversity\ PAU · #1.}\ }
\newcommand{\res}[1]{#1}\newcommand{\rest}[1]{#1}

% =====================================================================
% Capçalera FentMates de les fitxes: un color per a cada curs.
% Els colors són els de la pàgina inicial (LEVELS a assets/js/materials.js);
% el logo ve de _eines/logo/pdf/ (python3 _eines/logo/genera.py).
%
% Ús, al final del preàmbul de la unitat:
%   % =====================================================================
% Capçalera FentMates de les fitxes: un color per a cada curs.
% Els colors són els de la pàgina inicial (LEVELS a assets/js/materials.js);
% el logo ve de _eines/logo/pdf/ (python3 _eines/logo/genera.py).
%
% Ús, al final del preàmbul de la unitat:
%   % =====================================================================
% Capçalera FentMates de les fitxes: un color per a cada curs.
% Els colors són els de la pàgina inicial (LEVELS a assets/js/materials.js);
% el logo ve de _eines/logo/pdf/ (python3 _eines/logo/genera.py).
%
% Ús, al final del preàmbul de la unitat:
%   \input{../comu/fentmates.tex}
%   \fmcurs{1eso}
% i, al començament del document:
%   \capcalera{Matemàtiques}{1r ESO}{Fitxa d'exercicis · Unitat 2}{Divisibilitat}
% Opcional abans de \input: \def\fmneed{14} (espai mínim abans de cada secció).
% =====================================================================
\makeatletter
\def\fm@logodir{../../logo/pdf}
\providecommand{\fmneed}{10}
\def\fm@cursnom{}

% clau del curs, color i color del text damunt del color
\newcommand{\fm@curs}[3]{\@namedef{fm@color@#1}{#2}\@namedef{fm@text@#1}{#3}}
\fm@curs{1eso}{E6B84A}{navy}
\fm@curs{2eso}{3B82F6}{white}
\fm@curs{3eso}{FF7A50}{white}
\fm@curs{4eso-a}{A78BFA}{white}
\fm@curs{4eso-b}{6D28D9}{white}
\fm@curs{1batx-mi}{9B1C1C}{white}
\fm@curs{1batx-ccss}{EF6B6B}{white}
\fm@curs{2batx-mii}{1E6B47}{white}
\fm@curs{2batx-ccss}{2BB3A3}{white}

\newcommand{\fmcurs}[1]{%
  \@ifundefined{fm@color@#1}{\PackageError{fentmates}{Curs desconegut: #1}{}}{}%
  \edef\fm@tmp{\noexpand\definecolor{curs}{HTML}{\@nameuse{fm@color@#1}}%
               \noexpand\colorlet{curstext}{\@nameuse{fm@text@#1}}}\fm@tmp
  \colorlet{cursfosc}{curs!72!black}%
  \def\fm@logo{\fm@logodir/vertical-#1.pdf}%
  \def\fm@marca{\fm@logodir/marca-#1.pdf}}

% ---------- capçalera de la primera pàgina ----------
% #1 matèria (petit), #2 curs, #3 tipus de document i unitat, #4 títol
\newcommand{\capcalera}[4]{%
  \gdef\fm@cursnom{#2}%
  \par\noindent
  \begin{tcolorbox}[colback=curs,colframe=curs,boxrule=0pt,arc=2.2mm,width=3.5cm,height=2.15cm,
      valign=center,left=3mm,right=2mm,top=1mm,bottom=1mm,nobeforeafter,box align=center]
    \color{curstext}{\footnotesize #1}\par\vspace{1pt}{\fontsize{19}{21}\selectfont\bfseries #2\par}%
  \end{tcolorbox}\hspace{5mm}%
  \begin{minipage}[c]{\dimexpr\linewidth-3.5cm-5mm-2.9cm\relax}
    {\color{cursfosc}\small\bfseries\MakeUppercase{#3}}\par\vspace{2pt}
    {\color{navy}\fontsize{20}{23}\selectfont\bfseries #4\par}
  \end{minipage}\hfill
  \begin{minipage}[c]{2.6cm}\raggedleft\includegraphics[width=2.6cm]{\fm@logo}\end{minipage}\par
  \vspace{2pt}}

% ---------- barres de secció del color del curs ----------
\renewcommand{\seccio}[2]{\par\needspace{\fmneed\baselineskip}\bigskip
  \begin{tcolorbox}[colback=curs,colframe=curs,coltext=curstext,arc=2mm,boxrule=0pt,left=3mm,right=3mm,top=1.2mm,bottom=1.2mm]
  {\large\bfseries #1\quad #2}
  \end{tcolorbox}\vspace{-2pt}}

% ---------- estrelles de dificultat del color del curs (a Batxillerat no se'n posen) ----------
\@ifundefined{estrelles}{}{%
  \renewcommand{\estrelles}[1]{{\color{curs}\ifcase#1\or\faStar\faStar[regular]\faStar[regular]\or\faStar\faStar\faStar[regular]\or\faStar\faStar\faStar\fi}}}

% ---------- capçalera i peu de les pàgines ----------
% primera pàgina: logos institucionals al peu
\fancypagestyle{primera}{\fancyhf{}\renewcommand{\headrulewidth}{0pt}
  \fancyfoot[L]{\raisebox{-0.32\height}{\includegraphics[height=0.78cm]{logo-gva.png}}\hspace{5mm}%
    \raisebox{-0.32\height}{\includegraphics[height=0.78cm]{logo-vermellar.png}}}
  \fancyfoot[R]{\scriptsize\color{muted}\href{https://creativecommons.org/licenses/by-nc-sa/4.0/deed.ca}{\ccbyncsa}\ \ CC BY-NC-SA 4.0 · fentmates.org\qquad\small\thepage}}
% la resta: el símbol a l'esquerra i el curs a la dreta, amb el títol al centre com abans
\fancyhead[L]{\raisebox{-0.3\height}{\includegraphics[height=0.62cm]{\fm@marca}}}
\fancyhead[R]{\small\bfseries\color{cursfosc}\fm@cursnom}
\renewcommand{\headrule}{\color{curs!45}\hrule height \headrulewidth}
\makeatother

%   \fmcurs{1eso}
% i, al començament del document:
%   \capcalera{Matemàtiques}{1r ESO}{Fitxa d'exercicis · Unitat 2}{Divisibilitat}
% Opcional abans de \input: \def\fmneed{14} (espai mínim abans de cada secció).
% =====================================================================
\makeatletter
\def\fm@logodir{../../logo/pdf}
\providecommand{\fmneed}{10}
\def\fm@cursnom{}

% clau del curs, color i color del text damunt del color
\newcommand{\fm@curs}[3]{\@namedef{fm@color@#1}{#2}\@namedef{fm@text@#1}{#3}}
\fm@curs{1eso}{E6B84A}{navy}
\fm@curs{2eso}{3B82F6}{white}
\fm@curs{3eso}{FF7A50}{white}
\fm@curs{4eso-a}{A78BFA}{white}
\fm@curs{4eso-b}{6D28D9}{white}
\fm@curs{1batx-mi}{9B1C1C}{white}
\fm@curs{1batx-ccss}{EF6B6B}{white}
\fm@curs{2batx-mii}{1E6B47}{white}
\fm@curs{2batx-ccss}{2BB3A3}{white}

\newcommand{\fmcurs}[1]{%
  \@ifundefined{fm@color@#1}{\PackageError{fentmates}{Curs desconegut: #1}{}}{}%
  \edef\fm@tmp{\noexpand\definecolor{curs}{HTML}{\@nameuse{fm@color@#1}}%
               \noexpand\colorlet{curstext}{\@nameuse{fm@text@#1}}}\fm@tmp
  \colorlet{cursfosc}{curs!72!black}%
  \def\fm@logo{\fm@logodir/vertical-#1.pdf}%
  \def\fm@marca{\fm@logodir/marca-#1.pdf}}

% ---------- capçalera de la primera pàgina ----------
% #1 matèria (petit), #2 curs, #3 tipus de document i unitat, #4 títol
\newcommand{\capcalera}[4]{%
  \gdef\fm@cursnom{#2}%
  \par\noindent
  \begin{tcolorbox}[colback=curs,colframe=curs,boxrule=0pt,arc=2.2mm,width=3.5cm,height=2.15cm,
      valign=center,left=3mm,right=2mm,top=1mm,bottom=1mm,nobeforeafter,box align=center]
    \color{curstext}{\footnotesize #1}\par\vspace{1pt}{\fontsize{19}{21}\selectfont\bfseries #2\par}%
  \end{tcolorbox}\hspace{5mm}%
  \begin{minipage}[c]{\dimexpr\linewidth-3.5cm-5mm-2.9cm\relax}
    {\color{cursfosc}\small\bfseries\MakeUppercase{#3}}\par\vspace{2pt}
    {\color{navy}\fontsize{20}{23}\selectfont\bfseries #4\par}
  \end{minipage}\hfill
  \begin{minipage}[c]{2.6cm}\raggedleft\includegraphics[width=2.6cm]{\fm@logo}\end{minipage}\par
  \vspace{2pt}}

% ---------- barres de secció del color del curs ----------
\renewcommand{\seccio}[2]{\par\needspace{\fmneed\baselineskip}\bigskip
  \begin{tcolorbox}[colback=curs,colframe=curs,coltext=curstext,arc=2mm,boxrule=0pt,left=3mm,right=3mm,top=1.2mm,bottom=1.2mm]
  {\large\bfseries #1\quad #2}
  \end{tcolorbox}\vspace{-2pt}}

% ---------- estrelles de dificultat del color del curs (a Batxillerat no se'n posen) ----------
\@ifundefined{estrelles}{}{%
  \renewcommand{\estrelles}[1]{{\color{curs}\ifcase#1\or\faStar\faStar[regular]\faStar[regular]\or\faStar\faStar\faStar[regular]\or\faStar\faStar\faStar\fi}}}

% ---------- capçalera i peu de les pàgines ----------
% primera pàgina: logos institucionals al peu
\fancypagestyle{primera}{\fancyhf{}\renewcommand{\headrulewidth}{0pt}
  \fancyfoot[L]{\raisebox{-0.32\height}{\includegraphics[height=0.78cm]{logo-gva.png}}\hspace{5mm}%
    \raisebox{-0.32\height}{\includegraphics[height=0.78cm]{logo-vermellar.png}}}
  \fancyfoot[R]{\scriptsize\color{muted}\href{https://creativecommons.org/licenses/by-nc-sa/4.0/deed.ca}{\ccbyncsa}\ \ CC BY-NC-SA 4.0 · fentmates.org\qquad\small\thepage}}
% la resta: el símbol a l'esquerra i el curs a la dreta, amb el títol al centre com abans
\fancyhead[L]{\raisebox{-0.3\height}{\includegraphics[height=0.62cm]{\fm@marca}}}
\fancyhead[R]{\small\bfseries\color{cursfosc}\fm@cursnom}
\renewcommand{\headrule}{\color{curs!45}\hrule height \headrulewidth}
\makeatother

%   \fmcurs{1eso}
% i, al començament del document:
%   \capcalera{Matemàtiques}{1r ESO}{Fitxa d'exercicis · Unitat 2}{Divisibilitat}
% Opcional abans de \input: \def\fmneed{14} (espai mínim abans de cada secció).
% =====================================================================
\makeatletter
\def\fm@logodir{../../logo/pdf}
\providecommand{\fmneed}{10}
\def\fm@cursnom{}

% clau del curs, color i color del text damunt del color
\newcommand{\fm@curs}[3]{\@namedef{fm@color@#1}{#2}\@namedef{fm@text@#1}{#3}}
\fm@curs{1eso}{E6B84A}{navy}
\fm@curs{2eso}{3B82F6}{white}
\fm@curs{3eso}{FF7A50}{white}
\fm@curs{4eso-a}{A78BFA}{white}
\fm@curs{4eso-b}{6D28D9}{white}
\fm@curs{1batx-mi}{9B1C1C}{white}
\fm@curs{1batx-ccss}{EF6B6B}{white}
\fm@curs{2batx-mii}{1E6B47}{white}
\fm@curs{2batx-ccss}{2BB3A3}{white}

\newcommand{\fmcurs}[1]{%
  \@ifundefined{fm@color@#1}{\PackageError{fentmates}{Curs desconegut: #1}{}}{}%
  \edef\fm@tmp{\noexpand\definecolor{curs}{HTML}{\@nameuse{fm@color@#1}}%
               \noexpand\colorlet{curstext}{\@nameuse{fm@text@#1}}}\fm@tmp
  \colorlet{cursfosc}{curs!72!black}%
  \def\fm@logo{\fm@logodir/vertical-#1.pdf}%
  \def\fm@marca{\fm@logodir/marca-#1.pdf}}

% ---------- capçalera de la primera pàgina ----------
% #1 matèria (petit), #2 curs, #3 tipus de document i unitat, #4 títol
\newcommand{\capcalera}[4]{%
  \gdef\fm@cursnom{#2}%
  \par\noindent
  \begin{tcolorbox}[colback=curs,colframe=curs,boxrule=0pt,arc=2.2mm,width=3.5cm,height=2.15cm,
      valign=center,left=3mm,right=2mm,top=1mm,bottom=1mm,nobeforeafter,box align=center]
    \color{curstext}{\footnotesize #1}\par\vspace{1pt}{\fontsize{19}{21}\selectfont\bfseries #2\par}%
  \end{tcolorbox}\hspace{5mm}%
  \begin{minipage}[c]{\dimexpr\linewidth-3.5cm-5mm-2.9cm\relax}
    {\color{cursfosc}\small\bfseries\MakeUppercase{#3}}\par\vspace{2pt}
    {\color{navy}\fontsize{20}{23}\selectfont\bfseries #4\par}
  \end{minipage}\hfill
  \begin{minipage}[c]{2.6cm}\raggedleft\includegraphics[width=2.6cm]{\fm@logo}\end{minipage}\par
  \vspace{2pt}}

% ---------- barres de secció del color del curs ----------
\renewcommand{\seccio}[2]{\par\needspace{\fmneed\baselineskip}\bigskip
  \begin{tcolorbox}[colback=curs,colframe=curs,coltext=curstext,arc=2mm,boxrule=0pt,left=3mm,right=3mm,top=1.2mm,bottom=1.2mm]
  {\large\bfseries #1\quad #2}
  \end{tcolorbox}\vspace{-2pt}}

% ---------- estrelles de dificultat del color del curs (a Batxillerat no se'n posen) ----------
\@ifundefined{estrelles}{}{%
  \renewcommand{\estrelles}[1]{{\color{curs}\ifcase#1\or\faStar\faStar[regular]\faStar[regular]\or\faStar\faStar\faStar[regular]\or\faStar\faStar\faStar\fi}}}

% ---------- capçalera i peu de les pàgines ----------
% primera pàgina: logos institucionals al peu
\fancypagestyle{primera}{\fancyhf{}\renewcommand{\headrulewidth}{0pt}
  \fancyfoot[L]{\raisebox{-0.32\height}{\includegraphics[height=0.78cm]{logo-gva.png}}\hspace{5mm}%
    \raisebox{-0.32\height}{\includegraphics[height=0.78cm]{logo-vermellar.png}}}
  \fancyfoot[R]{\scriptsize\color{muted}\href{https://creativecommons.org/licenses/by-nc-sa/4.0/deed.ca}{\ccbyncsa}\ \ CC BY-NC-SA 4.0 · fentmates.org\qquad\small\thepage}}
% la resta: el símbol a l'esquerra i el curs a la dreta, amb el títol al centre com abans
\fancyhead[L]{\raisebox{-0.3\height}{\includegraphics[height=0.62cm]{\fm@marca}}}
\fancyhead[R]{\small\bfseries\color{cursfosc}\fm@cursnom}
\renewcommand{\headrule}{\color{curs!45}\hrule height \headrulewidth}
\makeatother

\fmcurs{2batx-mii}

\begin{document}
\thispagestyle{primera}

\capcalera{Matemàtiques II}{2n Batx}{Fitxa d'exercicis · Unitat 5}{\raggedright Límits i continuïtat}
\vspace{10pt}
Nom i cognoms: \rule{8.4cm}{0.4pt}\hfill Grup: \rule{1.6cm}{0.4pt}\hfill Data: \rule{2.2cm}{0.4pt}
\vspace{8pt}

\begin{tcolorbox}[colback=soft,colframe=soft,arc=2mm,left=3mm,right=3mm,top=1.5mm,bottom=1.5mm]
\small Resol els exercicis al quadern justificant cada pas. Cap límit d'aquesta fitxa necessita derivades (la regla de L'Hôpital la veurem més endavant). Els exercicis marcats amb {\color{teal}\bfseries\faUniversity\ PAU} són de les proves d'accés a la universitat d'altres comunitats autònomes, traduïts al valencià (només els apartats que corresponen a aquesta unitat).
\end{tcolorbox}

% =====================================================
\seccio{1}{Límits a l'infinit}

\ex{2} Calcula els límits següents:\begin{tasks}(3)
  \task $\displaystyle\lim_{x\to -\infty}(2x^5-x^2)$
  \task $\displaystyle\lim_{x\to +\infty}\dfrac{4x^3-x}{2x^3+7}$
  \task $\displaystyle\lim_{x\to -\infty}\dfrac{1-3x^2}{x+2}$
  \task $\displaystyle\lim_{x\to +\infty}\dfrac{x^2+1}{x^3-2}$
  \task $\displaystyle\lim_{x\to +\infty}\dfrac{2^{x+1}+5}{2^{x}-3}$
  \task $\displaystyle\lim_{x\to -\infty}\dfrac{3e^{2x}-e^{x}}{e^{2x}+4}$
\end{tasks}\recorda{Quan $x\to\pm\infty$, en un quocient de polinomis manen els termes de major grau. Amb exponencials, dividim per l'exponencial que creix més de pressa (i recorda que $e^{x}\to0$ quan $x\to-\infty$).}

\ex{2} Calcula els límits següents amb arrels. Compte amb el signe quan $x\to-\infty$:\begin{tasks}(3)
  \task $\displaystyle\lim_{x\to +\infty}\dfrac{\sqrt{9x^2+1}}{2x-5}$
  \task $\displaystyle\lim_{x\to -\infty}\dfrac{\sqrt{9x^2+1}}{2x-5}$
  \task $\displaystyle\lim_{x\to -\infty}\dfrac{3x+\sqrt{x^2+5}}{2x+1}$
  \task $\displaystyle\lim_{x\to -\infty}\dfrac{\sqrt[3]{8x^3-1}}{x}$
  \task $\displaystyle\lim_{x\to +\infty}\dfrac{\sqrt{x^3+x}}{x+4}$
  \task $\displaystyle\lim_{x\to -\infty}\dfrac{\sqrt{x^2-5x+3}}{3x-2}$
\end{tasks}\recorda{$\sqrt{x^2}=|x|$: quan $x\to-\infty$, $\sqrt{a x^2+\dots}\approx-\sqrt a\,x$. Si dubtes, fes el canvi $x=-t$ amb $t\to+\infty$.}

\ex{2} Ordena de menor a major ordre d'infinit (quan $x\to+\infty$) \ $x^3,\ 2^x,\ \ln x,\ \sqrt x,\ e^x,\ x^{10},\ (\ln x)^3$, \ i calcula:\begin{tasks}(3)
  \task $\displaystyle\lim_{x\to +\infty}\dfrac{x^3}{e^x}$
  \task $\displaystyle\lim_{x\to +\infty}\dfrac{\ln x}{x^2}$
  \task $\displaystyle\lim_{x\to +\infty}\dfrac{2^x+x^4}{3\cdot2^x}$
  \task $\displaystyle\lim_{x\to +\infty}\left(x^{10}-1{,}5^x\right)$
  \task $\displaystyle\lim_{x\to +\infty}\dfrac{3^x+2^x}{3^x-x^2}$
  \task $\displaystyle\lim_{x\to +\infty}\dfrac{(\ln x)^3}{\sqrt x}$
\end{tasks}

\ex{2} Calcula, fent servir les propietats dels logaritmes:\begin{tasks}(2)
  \task $\displaystyle\lim_{x\to +\infty}\big(\ln(2x+1)-\ln(x+3)\big)$
  \task $\displaystyle\lim_{x\to +\infty}\big(\ln(x^2+1)-\ln x\big)$
\end{tasks}

\ex{2} \pau{Andalusia, ordinària 2023} Considera la funció $f(x)=\dfrac{1}{e^x+e^{-x}}$. Calcula $\displaystyle\lim_{x\to+\infty}x^2f(x)$.


% =====================================================
\seccio{2}{Indeterminacions a l'infinit}

\ex{2} Indeterminació $\infty-\infty$:\begin{tasks}(3)
  \task $\displaystyle\lim_{x\to +\infty}\left(\dfrac{x^2}{x+1}-x\right)$
  \task $\displaystyle\lim_{x\to +\infty}\left(\dfrac{x^3}{x^2+1}-\dfrac{x^2}{x+1}\right)$
  \task $\displaystyle\lim_{x\to +\infty}\left(\sqrt{x^2+8x}-x\right)$
  \task $\displaystyle\lim_{x\to +\infty}\left(\sqrt{x^2+3x}-\sqrt{x^2-x}\right)$
  \task $\displaystyle\lim_{x\to -\infty}\left(\sqrt{x^2+2x}+x\right)$
  \task $\displaystyle\lim_{x\to +\infty}\left(\sqrt{x^2+1}-3x\right)$
\end{tasks}\recorda{Amb fraccions, fem la resta. Amb arrels, multipliquem i dividim pel \textbf{conjugat}. Abans, comprova si els termes principals es cancel·len: si no ho fan, no hi ha indeterminació real.}

\ex{3} Calcula $\displaystyle\lim_{x\to+\infty}\left(\sqrt[3]{x^3+3x^2}-x\right)$. \ \textit{Indicació}: $A^3-B^3=(A-B)(A^2+AB+B^2)$.

\ex{2} Indeterminació $0\cdot\infty$:\begin{tasks}(2)
  \task $\displaystyle\lim_{x\to +\infty}x\left(\sqrt{x^2+4}-x\right)$
  \task $\displaystyle\lim_{x\to +\infty}(2x+1)\left(\sqrt{x+1}-\sqrt x\right)$
  \task $\displaystyle\lim_{x\to +\infty}x\ln\left(1+\dfrac2x\right)$
  \task $\displaystyle\lim_{x\to +\infty}x^3e^{-x}$
\end{tasks}\recorda{Transformem el producte en un quocient ($f\cdot g=\frac{f}{1/g}$) o operem. Un exponent es pot posar dins del logaritme: $x\ln A=\ln A^{x}$.}

\ex{2} Indeterminació $1^{\infty}$:\begin{tasks}(3)
  \task $\displaystyle\lim_{x\to +\infty}\left(\dfrac{x+2}{x-1}\right)^{2x}$
  \task $\displaystyle\lim_{x\to +\infty}\left(\dfrac{2x-1}{2x+3}\right)^{x+1}$
  \task $\displaystyle\lim_{x\to +\infty}\left(\dfrac{x^2+3}{x^2-1}\right)^{x^2}$
  \task $\displaystyle\lim_{x\to +\infty}\left(\dfrac{3x+1}{x+2}\right)^{x}$
  \task $\displaystyle\lim_{x\to -\infty}\left(\dfrac{x-4}{x}\right)^{3x}$
  \task $\displaystyle\lim_{x\to +\infty}\left(\dfrac{x^2-x+1}{x^2+2}\right)^{\frac{x^2+1}{x-3}}$
\end{tasks}\recorda{Si $f\to1$ i $g\to\pm\infty$: \ $\lim f^{\,g}=e^{\lim\,(f-1)\,g}$. Si la base no tendeix a 1, no hi ha indeterminació.}

\ex{3} Indeterminacions $\infty^0$. Calcula prenent logaritmes:\begin{tasks}(2)
  \task $\displaystyle\lim_{x\to +\infty}x^{1/x}$
  \task $\displaystyle\lim_{x\to +\infty}(3x)^{\frac{2}{\ln x}}$
\end{tasks}

\ex{2} \pau{Aragó, ordinària 2024} Calcula justificadament $\displaystyle\lim_{x\to+\infty}\left[\sqrt{x^2+5}-(x+2)\right]$.

\ex{2} \pau{Aragó, ordinària 2022} Calcula $\displaystyle\lim_{x\to+\infty}\left[\sqrt{3x^2-2}-\left(\sqrt3\,x+5\right)\right]$.

\ex{2} \pau{Navarra, extraordinària 2021} Calcula $\displaystyle\lim_{x\to+\infty}\dfrac{\sqrt x}{\sqrt{3x^3+2x^2}-\sqrt{3x^3}}$.

\ex{2} \pau{Múrcia, ordinària 2024} Calcula:\begin{tasks}(2)
  \task $\displaystyle\lim_{x\to +\infty}\left(\sqrt{x+9}-\sqrt{x-9}\right)$
  \task $\displaystyle\lim_{x\to +\infty}\dfrac{\ln x}{\sqrt x}$
\end{tasks}

\ex{2} \pau{Castella-la Manxa, ordinària 2023} Calcula $\displaystyle\lim_{x\to+\infty}\left(\dfrac{5x+1}{5x}\right)^{x^2}$.

\ex{2} \pau{Navarra, extraordinària 2022} Calcula $\displaystyle\lim_{x\to+\infty}\left(\dfrac{x^2}{x^2+x-1}\right)^{2x-1}$.

\ex{2} \pau{La Rioja, extraordinària 2023} Calcula $\displaystyle\lim_{x\to+\infty}(1+x)^{1/x}$.

\ex{2} \pau{Illes Balears, extraordinària 2021} El nombre d'individus, en milions, d'una població és $P(t)=\dfrac{15+t^2}{(t+1)^2}$, on $t\geq0$ són els anys transcorreguts des de l'1 de gener del 2000.\begin{tasks}(1)
  \task Quina població hi havia l'1 de gener del 2000?
  \task Calcula la mida de la població a llarg termini.
\end{tasks}


% =====================================================
\seccio{3}{Límits en un nombre i indeterminacions}

\ex{2} Calcula els límits laterals i digues si existeix el límit:\begin{tasks}(2)
  \task $\displaystyle\lim_{x\to 2}\dfrac{x+1}{(x-2)^2}$
  \task $\displaystyle\lim_{x\to 3}\dfrac{x-3}{x^3-5x^2+3x+9}$
\end{tasks}

\ex{2} Indeterminació $\frac00$ amb polinomis:\begin{tasks}(3)
  \task $\displaystyle\lim_{x\to 2}\dfrac{x^3-8}{x-2}$
  \task $\displaystyle\lim_{x\to 1}\dfrac{x^4-1}{x^2-x}$
  \task $\displaystyle\lim_{x\to 2}\dfrac{x^4-4x^3+16x-16}{x^3-6x^2+12x-8}$
\end{tasks}\recorda{Si $P(a)=Q(a)=0$, tots dos són divisibles per $(x-a)$: factoritzem (Ruffini), simplifiquem i tornem a substituir.}

\ex{2} Indeterminació $\frac00$ amb arrels:\begin{tasks}(3)
  \task $\displaystyle\lim_{x\to 4}\dfrac{\sqrt{x+5}-3}{x-4}$
  \task $\displaystyle\lim_{x\to 3}\dfrac{\sqrt{x+1}-\sqrt{2x-2}}{x^2-9}$
  \task $\displaystyle\lim_{x\to 0}\dfrac{\sqrt{1+x}-1}{\sqrt{4+x}-2}$
\end{tasks}

\ex{2} Indeterminacions $\infty-\infty$ i $1^\infty$ en un nombre:\begin{tasks}(3)
  \task $\displaystyle\lim_{x\to 1}\left(\dfrac{1}{x-1}-\dfrac{2}{x^2-1}\right)$
  \task $\displaystyle\lim_{x\to 2}\left(\dfrac{x}{x-2}-\dfrac{8}{x^2-4}\right)$
  \task $\displaystyle\lim_{x\to 0}(1+3x)^{2/x}$
  \task $\displaystyle\lim_{x\to 2}\left(\dfrac{x+2}{2x}\right)^{\frac{1}{x-2}}$
  \task $\displaystyle\lim_{x\to 1}x^{\frac{3}{x-1}}$
  \task $\displaystyle\lim_{x\to 1}(2x-1)^{\frac{1}{x-1}}$
\end{tasks}

\ex{2} \pau{Castella-la Manxa, extraordinària 2024} Calcula $\displaystyle\lim_{x\to-2}\dfrac{x^3+2x^2+x+2}{x^2+2x}$.

\ex{2} \pau{Castella-la Manxa, extraordinària 2023} Calcula $\displaystyle\lim_{x\to3}\dfrac{x^3-3x^2+3x-9}{3x-9}$.

\ex{2} \pau{Castella i Lleó, ordinària 2023} Calcula $\displaystyle\lim_{x\to2}\dfrac{\sqrt{x^3+x-1}-\sqrt{x^3+1}}{x-2}$.

\ex{2} \pau{Múrcia, ordinària 2023} Calcula $\displaystyle\lim_{x\to0}\dfrac{\sqrt{9+x}-\sqrt{9-x}}{3x}$.


% =====================================================
\seccio{4}{Límits trigonomètrics i infinitèsims equivalents}

\ex{2} Calcula, fent servir infinitèsims equivalents:\begin{tasks}(3)
  \task $\displaystyle\lim_{x\to 0}\dfrac{\sin 4x}{3x}$
  \task $\displaystyle\lim_{x\to 0}\dfrac{x\tan 2x}{1-\cos x}$
  \task $\displaystyle\lim_{x\to 0}\dfrac{\ln(1+x^2)}{\sin^2 3x}$
  \task $\displaystyle\lim_{x\to 0}\dfrac{\sin 3x}{\tan 5x}$
  \task $\displaystyle\lim_{x\to 0}\dfrac{1-\cos 2x}{x\sin x}$
  \task $\displaystyle\lim_{x\to 0}\dfrac{e^{2x}-1}{\ln(1+3x)}$
\end{tasks}\recorda{Si $u\to0$: \ $\sin u\sim u$, \ $\tan u\sim u$, \ $1-\cos u\sim\frac{u^2}{2}$, \ $e^u-1\sim u$, \ $\ln(1+u)\sim u$. Només es poden substituir en \textbf{productes i quocients}, mai en sumes o restes.}

\ex{2} Calcula (si el punt no és 0, fes primer un canvi de variable):\begin{tasks}(3)
  \task $\displaystyle\lim_{x\to 0}\dfrac{e^{x}-e^{-x}}{\sin x}$
  \task $\displaystyle\lim_{x\to \frac\pi2}\dfrac{\cos x}{2x-\pi}$
  \task $\displaystyle\lim_{x\to \pi}\dfrac{\sin x}{x-\pi}$
  \task $\displaystyle\lim_{x\to 1}\dfrac{\ln x}{x^2-1}$
  \task $\displaystyle\lim_{x\to 0}\dfrac{\sin x-\tan x}{x^2\sin x}$
  \task $\displaystyle\lim_{x\to 0}(\cos x)^{1/x^2}$
\end{tasks}

\ex{2} Calcula, fent servir que el producte d'una funció fitada per una que tendeix a 0 tendeix a 0:\begin{tasks}(3)
  \task $\displaystyle\lim_{x\to -\infty}\dfrac{3x+\cos x}{x}$
  \task $\displaystyle\lim_{x\to 0}x^2\cos\dfrac1x$
  \task $\displaystyle\lim_{x\to +\infty}x\sin\dfrac1x$
\end{tasks}

\ex{2} \pau{Múrcia, ordinària 2023} Calcula $\displaystyle\lim_{x\to0}\dfrac{\cos(2x)-1}{x\sin x}$.

\ex{2} \pau{Castella i Lleó, extraordinària 2022} Calcula $\displaystyle\lim_{x\to0}\dfrac{\ln(1+x)}{e^x-1}$.

\ex{2} \pau{Galícia, ordinària 2026} Calcula els límits següents:\begin{tasks}(2)
  \task $\displaystyle\lim_{x\to 1}\dfrac{\ln x}{x}$
  \task $\displaystyle\lim_{x\to 1}\dfrac{x-1}{\ln x}$
  \task $\displaystyle\lim_{x\to 0}\dfrac{\tan x}{\sin x}$
  \task $\displaystyle\lim_{x\to 0^+}\sqrt x\,\ln x$
\end{tasks}

\ex{2} \pau{La Rioja, extraordinària 2023} Calcula $\displaystyle\lim_{x\to0}(\cos2x)^{3/x^2}$.

\ex{3} \pau{Aragó, extraordinària 2021} Calcula $\displaystyle\lim_{x\to1}\left(\sin\dfrac{\pi x}{2}\right)^{\frac{1}{(1-x)^2}}$.

\ex{2} \pau{Navarra, extraordinària 2021} Calcula $\displaystyle\lim_{x\to\infty}x\sin\dfrac1x$.

\ex{2} \pau{Múrcia, extraordinària 2025} Calcula $\displaystyle\lim_{x\to+\infty}x\left(e^{1/x}-1\right)$.


% =====================================================
\seccio{5}{Límits amb paràmetres}

\ex{2} Calcula el valor dels paràmetres perquè:\begin{tasks}(1)
  \task $\displaystyle\lim_{x\to +\infty}\dfrac{ax^2-3x}{2x^2+5}=4$
  \task $\displaystyle\lim_{x\to +\infty}\dfrac{ax^3+bx^2+1}{2x^2-x}=3$
  \task $\displaystyle\lim_{x\to +\infty}\left(\dfrac{2x+3}{2x+a}\right)^{4x}=e^2$
  \task $\displaystyle\lim_{x\to 0}\dfrac{e^{mx}-1}{\sin 2x}=\dfrac32$
\end{tasks}\recorda{Calculem el límit com si el paràmetre fóra un nombre i imposem la condició.}

\ex{2} Troba el valor de $a$ perquè el límit siga finit, i calcula'l en cada cas:\begin{tasks}(2)
  \task $\displaystyle\lim_{x\to +\infty}\left(\sqrt{ax^2+x}-2x\right)$
  \task $\displaystyle\lim_{x\to -3}\dfrac{x^2+ax+3}{x^2+3x}$
\end{tasks}

\ex{2} Troba $a$ i $b$ perquè $\displaystyle\lim_{x\to1}\dfrac{x^2+ax+b}{x-1}=3$.

\ex{2} Troba $a$ perquè:\begin{tasks}(2)
  \task $\displaystyle\lim_{x\to +\infty}\left(\sqrt{x^2+ax}-\sqrt{x^2-x}\right)=2$
  \task $\displaystyle\lim_{x\to 0}\dfrac{1-\cos(ax)}{x\sin x}=8$
\end{tasks}

\ex{3} Discuteix, segons els valors del paràmetre $a$, els límits següents:\begin{tasks}(1)
  \task $\displaystyle\lim_{x\to +\infty}\dfrac{(a-1)x^2+x}{3x+1}$
  \task $\displaystyle\lim_{x\to +\infty}\left(\dfrac{x+1}{x+2}\right)^{ax^2+x}$
\end{tasks}

\ex{2} \pau{Andalusia, extraordinària 2022} Calcula $a$ sabent que $\displaystyle\lim_{x\to+\infty}\dfrac{ax}{(\ln x)^2+2x}=1$.

\ex{3} \pau{Aragó, extraordinària 2025} Calcula (si existeix), en funció del valor de $k\in\mathbb Z$, el límit \ $\displaystyle\lim_{x\to1}\dfrac{x^4+x^3-x^2-x}{\left(x^2-1\right)^{2k}}$.

\ex{2} \pau{Castella-la Manxa, extraordinària 2022} Troba raonadament $a$ i $b$ perquè la funció $f(x)=\dfrac{ax+1}{2x+b}$ tinga una discontinuïtat de salt infinit en $x=1$ i tendisca a 2 quan $x\to+\infty$.


% =====================================================
\seccio{6}{Continuïtat de funcions}

\ex{2} Estudia la continuïtat en $x=0$ de les funcions següents i classifica la discontinuïtat, si n'hi ha:\begin{tasks}(1)
  \task $f(x)=\begin{cases}\dfrac{\sin 3x}{x}&\text{si }x\neq0\\[3pt]3&\text{si }x=0\end{cases}$
  \task $g(x)=\dfrac{1}{1+e^{1/x}}$
  \task $h(x)=x\cos\dfrac1x$. Quin valor caldria donar a $h(0)$ perquè fóra contínua?
\end{tasks}\recorda{$f$ és contínua en $a$ si $\exists f(a)$, $\exists\lim_{x\to a}f(x)$ i són iguals. Discontinuïtats: evitable (el límit existeix), de salt finit (laterals finits i diferents), de salt infinit (un lateral infinit) i asimptòtica (tots dos infinits).}

\ex{2} Estudia la continuïtat de les funcions següents:\begin{tasks}(1)
  \task $f(x)=\begin{cases}\cos x&\text{si }x\leq0\\[3pt]\dfrac{e^x-1}{x}&\text{si }0<x<2\\[3pt]\log_2(x+2)&\text{si }x\geq2\end{cases}$
  \task $g(x)=\begin{cases}\dfrac{3}{x+2}&\text{si }x<-1\\[3pt]2-|x|&\text{si }-1\leq x\leq1\\[3pt]\sqrt{x+3}-1&\text{si }x>1\end{cases}$
  \task $h(x)=\begin{cases}\dfrac{\sin 2x}{x}&\text{si }x<0\\[3pt]x+2&\text{si }0\leq x<\pi\\[3pt]\cos x+\pi+3&\text{si }x\geq\pi\end{cases}$
\end{tasks}\recorda{Primer, cada tros en el seu interval obert (denominadors, logaritmes, arrels, $\tan$\ldots\ només compten si el punt és dins del tros). Després, els punts on canvia l'expressió.}

\ex{2} Estudia la continuïtat de:\begin{tasks}(3)
  \task $f(x)=\dfrac{x}{1+e^{1/x}}$
  \task $g(x)=\dfrac{x^2-4}{|x-2|}$
  \task $h(x)=e^{\frac{1}{x-1}}$
\end{tasks}

\ex{2} \pau{Castella-la Manxa, extraordinària 2024} Estudia la continuïtat de la funció i, si en té, classifica'n les discontinuïtats:\[f(x)=\begin{cases}x^2+2x+1&\text{si }x<3\\[3pt]\dfrac{2x}{x-4}&\text{si }x\geq3\end{cases}\]

\ex{2} \pau{Castella-la Manxa, extraordinària 2020} Estudia la continuïtat de $f$ en $x=1$ i en $x=2$ i classifica-hi el tipus de discontinuïtat, si n'hi ha:\[f(x)=\begin{cases}2^{x-1}&\text{si }x\leq1\\[3pt]x-2&\text{si }1<x<2\\[3pt]\ln(x-1)&\text{si }x\geq2\end{cases}\]

\ex{2} \pau{Astúries, ordinària 2025} Estudia si la funció és contínua en el seu domini:\[f(x)=\begin{cases}\dfrac x2&\text{si }0\leq x<4\\[3pt]3-(x-5)^2&\text{si }4\leq x\end{cases}\]

\ex{2} \pau{Madrid, ordinària 2023} Donada la funció\[f(x)=\begin{cases}\dfrac{x^2}{2+x^2}&\text{si }x\leq-1\\[3pt]\dfrac{2x^2}{3-3x}&\text{si }x>-1\end{cases}\]\begin{tasks}(1)
  \task Estudia la continuïtat de $f$ en $\mathbb R$.
  \task Calcula $\displaystyle\lim_{x\to-\infty}f(x)^{2x^2-1}$.
\end{tasks}

\ex{2} \pau{Illes Balears, extraordinària 2024} Indica el domini de $g(x)=\dfrac{(x-1)\sqrt x}{x^2-1}$. És contínua en el seu domini? Si no ho és, quin tipus de discontinuïtat presenta?

\ex{2} \pau{Castella i Lleó, extraordinària 2021} Estudia la continuïtat de la funció\[f(x)=\begin{cases}\dfrac{1-\cos x}{x}&\text{si }x\neq0\\[3pt]0&\text{si }x=0\end{cases}\]

\ex{2} \pau{Galícia, extraordinària 2025} Estudia la continuïtat en $x=0$ de la funció\[f(x)=\begin{cases}x\,e^{4x}&\text{si }x<0\\[3pt]\dfrac{\ln(1+x)}{1+x}&\text{si }x\geq0\end{cases}\]

\ex{3} \pau{Navarra, ordinària 2023} Estudia la continuïtat en $\mathbb R$ de la funció\[f(x)=\begin{cases}\dfrac{\cos^2(\pi x)-1}{1-x}&\text{si }x<1\\[3pt]\ln\left(x\,e^{x+1}\right)-2x&\text{si }x\geq1\end{cases}\]

\ex{2} \pau{Madrid, extraordinària 2022} Estudia la continuïtat en $x=0$ de $f(x)=\begin{cases}x&\text{si }x\leq0\\x\ln x&\text{si }x>0\end{cases}$.


% =====================================================
\seccio{7}{Funcions a trossos amb paràmetres}

\ex{2} Troba els paràmetres perquè les funcions siguen contínues en tot el seu domini:\begin{tasks}(2)
  \task $f(x)=\begin{cases}\dfrac{\sin(ax)}{x}&\text{si }x<0\\[3pt]x^2+3&\text{si }x\geq0\end{cases}$
  \task $g(x)=\begin{cases}\dfrac{e^{ax}-1}{x}&\text{si }x<0\\[3pt]\cos x+3&\text{si }x\geq0\end{cases}$
  \task $h(x)=\begin{cases}\dfrac{x^2+ax-10}{x-2}&\text{si }x<2\\[3pt]b+\ln(x-1)&\text{si }x\geq2\end{cases}$
  \task $k(x)=\begin{cases}\dfrac{\tan3x}{\sin x}&\text{si }-\frac\pi6<x<0\\[3pt]a&\text{si }x=0\\[3pt]\sqrt{x+b}&\text{si }x>0\end{cases}$
\end{tasks}\recorda{En cada punt de canvi igualem els límits laterals i el valor de la funció. Si un tros dona $\frac00$, resolem la indeterminació; si el límit ha de ser finit i el denominador s'anul·la, el numerador també s'ha d'anul·lar.}

\ex{2} Troba els paràmetres perquè les funcions siguen contínues en $\mathbb R$:\begin{tasks}(1)
  \task $f(x)=\begin{cases}-2\sin x&\text{si }x\leq-\frac\pi2\\[3pt]a\sin x+b&\text{si }-\frac\pi2<x<\frac\pi2\\[3pt]\cos x&\text{si }x\geq\frac\pi2\end{cases}$
  \task $g(x)=\begin{cases}a\,e^{x}&\text{si }x\leq0\\[3pt]\dfrac{\ln(1+3x)}{x}&\text{si }0<x\leq1\\[3pt]b\sqrt{x+3}&\text{si }x>1\end{cases}$
\end{tasks}

\ex{2} Troba $k$ i $a$ perquè les funcions siguen contínues:\begin{tasks}(2)
  \task $f(x)=\begin{cases}\dfrac{\sqrt{x+1}-2}{x-3}&\text{si }x\neq3\\[3pt]k&\text{si }x=3\end{cases}$
  \task $g(x)=\begin{cases}x+e^2&\text{si }x\leq0\\[3pt](1+ax)^{1/x}&\text{si }x>0\end{cases}$
\end{tasks}

\ex{2} \pau{Andalusia, ordinària 2022} Considera la funció contínua\[f(x)=\begin{cases}\dfrac1x&\text{si }x<-1\\[3pt]ax+b&\text{si }-1\leq x<1\\[3pt]\dfrac{x^2}{x+1}&\text{si }x\geq1\end{cases}\]\begin{tasks}(1)
  \task Calcula $a$ i $b$.
  \task Estudia i troba les asímptotes de la gràfica de $f$.
\end{tasks}

\ex{2} \pau{Aragó, extraordinària 2024} Estudia la continuïtat en $\mathbb R$, segons els valors de $a$ i $b$, de\[f(x)=\begin{cases}a-\cos x&\text{si }x\leq0\\[3pt]x^2-b\sin\left(x+\frac\pi2\right)&\text{si }x>0\end{cases}\]

\ex{2} \pau{Aragó, ordinària 2024} Estudia la continuïtat en $\mathbb R$, segons els valors de $a$, de\[f(x)=\begin{cases}\dfrac{e^{2x}-1}{x}&\text{si }x\neq0\\[3pt]a&\text{si }x=0\end{cases}\]

\ex{2} \pau{Aragó, ordinària 2022} Determina $a$ perquè $f$ siga contínua en $\mathbb R$:\[f(x)=\begin{cases}\sqrt{-x}+e^3&\text{si }x\leq0\\[3pt](1-x)^{a/x}&\text{si }x>0\end{cases}\]

\ex{2} \pau{Aragó, extraordinària 2021} Calcula els valors de $a\neq0$ perquè $f$ siga contínua:\[f(x)=\begin{cases}5-ax^2&\text{si }x\leq1\\[3pt]\dfrac{6}{ax}&\text{si }x>1\end{cases}\]

\ex{2} \pau{Illes Balears, ordinària 2022} Troba la condició que han de complir $a$ i $b$ perquè $f$ siga contínua:\[f(x)=\begin{cases}\dfrac{x^2+a}{2x-4}&\text{si }x\leq0\\[3pt]10x^2+x+b&\text{si }x>0\end{cases}\]

\ex{2} \pau{Illes Balears, extraordinària 2021} Considera $f(x)=\begin{cases}\dfrac{e^{ax}-1}{2x}&\text{si }x\neq0\\b&\text{si }x=0\end{cases}$.\begin{tasks}(1)
  \task Estudia la continuïtat de $f$ en els punts $x\neq0$.
  \task Quina relació hi ha d'haver entre $a$ i $b$ perquè $f$ siga contínua en $x=0$?
\end{tasks}

\ex{2} \pau{Canàries, ordinària 2021} Siga $f(x)=\dfrac{ax^2-2}{b-x}$.\begin{tasks}(1)
  \task Calcula $a$ i $b$ perquè $f(-2)=2$ i $f$ siga contínua en $\mathbb R-\{5\}$.
  \task Troba les asímptotes de $f$ si $a=-1$ i $b=-3$.
\end{tasks}

\ex{2} \pau{Cantàbria, extraordinària 2026} Considera\[f(x)=\begin{cases}x^2+3x+a^2&\text{si }x<1\\[3pt]\dfrac{\ln x}{x}+4a&\text{si }x\geq1\end{cases}\]\begin{tasks}(1)
  \task Calcula els valors de $a$ per als quals $f$ és contínua en tot el seu domini.
  \task Calcula les asímptotes de $f$ en funció de $a$.
\end{tasks}

\ex{2} \pau{Cantàbria, ordinària 2025} Considera $f(x)=\begin{cases}\dfrac{x^3-8}{x^2-4}&\text{si }x\neq2\\e^{k}&\text{si }x=2\end{cases}$.\begin{tasks}(1)
  \task Determina $k$ perquè $f$ siga contínua en $x=2$.
  \task Troba les asímptotes de $f$ i especifica'n el tipus.
\end{tasks}

\ex{2} \pau{Cantàbria, extraordinària 2020} Considera\[f(x)=\begin{cases}\sin x&\text{si }x\leq\frac\pi2\\[3pt]\dfrac2x+a&\text{si }x>\frac\pi2\end{cases}\]\begin{tasks}(1)
  \task Troba $a$ perquè $f$ siga contínua.
  \task Calcula $\displaystyle\lim_{x\to+\infty}f(x)$ i $\displaystyle\lim_{x\to-\infty}f(x)$.
\end{tasks}

\ex{2} \pau{Cantàbria, extraordinària 2018} Siga\[f(x)=\begin{cases}x+2&\text{si }x\leq-2\\[3pt]x^2+ax&\text{si }-2<x<0\\[3pt]2\sin x+b&\text{si }0\leq x\end{cases}\]\begin{tasks}(1)
  \task Determina $a$ i $b$ perquè $f$ siga contínua en $\mathbb R$.
  \task Si $a=3$ i $b=0$, classifica la discontinuïtat en $x=-2$.
\end{tasks}

\ex{2} \pau{Castella-la Manxa, extraordinària 2017} Determina $k$ perquè $f$ siga contínua en $x=0$:\[f(x)=\begin{cases}\left(\dfrac{x+1}{2x+1}\right)^{1/x}&\text{si }x<0\\[3pt]6x+k&\text{si }x\geq0\end{cases}\]

\ex{2} \pau{Castella i Lleó, extraordinària 2024} Quins valors han de prendre $a\in\mathbb R$ i $b\in\mathbb R-\{0\}$ perquè $f$ siga contínua en $\mathbb R$?\[f(x)=\begin{cases}\dfrac{\cos x-a}{bx^2}&\text{si }x\neq0\\[3pt]1&\text{si }x=0\end{cases}\]

\ex{2} \pau{Catalunya, ordinària 2026} Determina els valors de $m$ que fan que $f$ siga contínua en tot el seu domini:\[f(x)=\begin{cases}5e^{2x}&\text{si }x\leq0\\[3pt](x+m)^2+1&\text{si }0<x<2\\[3pt]1&\text{si }2\leq x\end{cases}\]

\ex{2} \pau{Extremadura, extraordinària 2023} Troba $a$ i $b$ perquè $f$ siga contínua en $x=1$ i la seua gràfica passe pel punt $(-1,5)$:\[f(x)=\begin{cases}2x^2+ax+b&\text{si }x\leq1\\[3pt]\ln x&\text{si }x>1\end{cases}\]

\ex{2} \pau{Múrcia, extraordinària 2022} Considera $f(x)=\begin{cases}\dfrac{\ln x}{x-1}&\text{si }x>0,\ x\neq1\\a&\text{si }x=1\end{cases}$.\begin{tasks}(1)
  \task Calcula el límit de $f(x)$ quan $x\to+\infty$.
  \task Determina $a$ perquè $f$ siga contínua en $x=1$.
\end{tasks}


% =====================================================
\seccio{8}{Teorema de Bolzano}

\ex{2} Demostra que l'equació $x^3+x-1=0$ té almenys una solució en l'interval $(0,1)$ i troba-la amb una xifra decimal exacta.\recorda{\textbf{Bolzano}: si $f$ és contínua en $[a,b]$ i $f(a)\cdot f(b)<0$, existeix almenys un $c\in(a,b)$ amb $f(c)=0$. Per afinar, parteix l'interval per la meitat i queda't amb la meitat on hi ha canvi de signe.}

\ex{2} Demostra que:\begin{tasks}(1)
  \task l'equació $x^5+x^3=1$ té alguna solució real, i troba un interval de longitud $0{,}1$ que la continga;
  \task l'equació $e^{-x}=x$ té alguna solució en $(0,1)$;
  \task les gràfiques de $y=\cos x$ i $y=x$ es tallen en algun punt de $\left(0,\frac\pi2\right)$;
  \task la funció $f(x)=x^3-x$ pren el valor $5$ en algun punt de $[1,2]$.
\end{tasks}

\ex{2} Raona si es pot aplicar el teorema de Bolzano i, si és el cas, què en podem deduir:\begin{tasks}(1)
  \task $f(x)=\dfrac1x$ en $[-1,1]$
  \task $f(x)=\tan x$ en $\left[\frac\pi4,\frac{3\pi}4\right]$
  \task $f(x)=x^3+ax-1$ en $[0,1]$, amb $a>0$
  \task $f(x)=x^2-1$ en $[-2,2]$
\end{tasks}

\ex{2} \pau{Illes Balears, ordinària 2025} L'altitud (en metres) d'un terreny en un tram de 500 metres és $A(x)=-0{,}0001x^3+0{,}05x^2-4x+200$, amb $x\in[0,500]$ la distància recorreguda horitzontalment. Demostra que hi ha almenys un punt del tram on l'altitud és 0.

\ex{2} \pau{Castella-la Manxa, ordinària 2023} Enuncia el teorema de Bolzano i utilitza'l per justificar que $f(x)=x^3+6x^2+3x-10$ té almenys una arrel en $[0,2]$.

\ex{2} \pau{Castella-la Manxa, extraordinària 2017} Enuncia el teorema de Bolzano i comprova si l'equació $\cos x=2-x$ té alguna solució real en l'interval $[0,2\pi]$.

\ex{2} \pau{Castella i Lleó, extraordinària 2022} Esbrina si la funció $f(x)=x+\sin x-2$ s'anul·la en algun punt de l'interval $\left[0,\frac\pi2\right]$.

\ex{2} \pau{Castella i Lleó, extraordinària 2024} Donada la funció $f(x)=e^x+x^3-2$, demostra que $f$ s'anul·la per a algun valor de $x$ i que aquest valor és únic.

\ex{2} \pau{Castella i Lleó, ordinària 2017} Enuncia el teorema de Bolzano, interpreta'l geomètricament i troba un interval en què $P(x)=x^6+x^4-1$ tinga almenys una arrel.

\ex{2} \pau{Extremadura, ordinària 2018} Enuncia el teorema de Bolzano i demostra, usant-lo, que la funció $f(x)=x^3+x-3$ té una arrel real positiva.

\ex{2} \pau{Madrid, ordinària 2022} Comprova si $f(x)=\dfrac{x}{x^2+1}$ verifica les hipòtesis del teorema de Bolzano en l'interval $[-1,1]$.

\ex{3} \pau{Catalunya, ordinària 2019} Considera la funció $f(x)=\dfrac{2x^3-5x+4}{1-x}$.\begin{tasks}(1)
  \task Calcula el domini i estudia la continuïtat de $f$. Té asímptota vertical?
  \task Observa que $f(-2)=-\frac23$, $f(0)=4$ i $f(2)=-10$. Raona si, a partir d'aquesta informació, podem deduir que l'interval $(-2,0)$ conté un zero de la funció. I l'interval $(0,2)$? Troba un interval determinat per dos enters consecutius que continga, com a mínim, un zero de la funció.
\end{tasks}

\ex{2} \pau{Galícia, ordinària 2026} Dibuixa, si és possible, la gràfica d'una funció contínua en $[a,b]$, amb $f(a)f(b)<0$, i que talle l'eix $X$ en més d'un punt de l'interval $(a,b)$.

\ex{2} \pau{Navarra, ordinària 2023} Es considera $f(x)=(x+1)\sin(\pi x)$. Comprova que existeix un valor $\alpha\in(0,1)$ tal que $f(\alpha)=\frac34$. Enuncia el resultat teòric utilitzat.

\ex{3} \pau{Navarra, extraordinària 2022} Siga $f(x)=\sin\left(\dfrac\pi4\ln\dfrac1x\right)$.\begin{tasks}(1)
  \task Demostra que és contínua en $\left[\frac1e,e\right]$.
  \task Demostra que existeix $\alpha\in\left(\frac1e,e\right)$ tal que $f(\alpha)=\dfrac{e\sqrt2}{1-e^2}$. Enuncia el resultat utilitzat.
\end{tasks}

\ex{3} \pau{Navarra, extraordinària 2021} Es considera $f(x)=\log_2\left[\sin\dfrac{\pi(x+1)}{4}+2^{\frac{x-5}{2}}\right]$.\begin{tasks}(1)
  \task Demostra que és contínua en $[6,7]$.
  \task Demostra que existeix $\alpha\in(6,7)$ tal que $f(\alpha)=0$.
\end{tasks}


% =====================================================
\seccio{9}{Asímptotes}

\ex{2} Calcula totes les asímptotes de les funcions següents i estudia els límits laterals en les verticals:\begin{tasks}(2)
  \task $f(x)=\dfrac{2x-1}{x+3}$
  \task $f(x)=\dfrac{2x^3-x}{x^2-1}$
  \task $f(x)=\dfrac{e^x}{x-1}$
  \task $f(x)=\sqrt{4x^2+x}$
  \task $f(x)=\ln\dfrac{x}{x-2}$
  \task $f(x)=x+\dfrac{\sin x}{x}$
\end{tasks}\recorda{AV: on la funció tendeix a $\pm\infty$ (zeros de denominadors, extrems del domini d'un logaritme\ldots). AH: $\lim_{x\to\pm\infty}f(x)=k$. AO $y=mx+n$: $m=\lim\frac{f(x)}{x}$, $n=\lim(f(x)-mx)$; estudia cada costat per separat.}

\ex{2} \pau{Andalusia, ordinària 2021} La gràfica de $f(x)=\dfrac{ax^2+bx+2}{x-1}$ ($x\neq1$) té una asímptota obliqua que passa pel punt $(1,1)$ i té pendent 2. Calcula $a$ i $b$.

\ex{2} \pau{Cantàbria, ordinària 2026} Calcula els valors de $a$ per als quals $f(x)=\dfrac{x^2+ax+1}{x+a^2}$ té com a asímptota obliqua la recta $y=x-6$.

\ex{2} \pau{Cantàbria, ordinària 2026} Calcula, si n'hi ha, les asímptotes de $f(x)=\dfrac{e^x-1}{e^{2x}-4}$.

\ex{2} \pau{Cantàbria, extraordinària 2018} Per a quin valor de $d$ té la funció $\dfrac{x^d+1}{x-2}$ una asímptota obliqua en $+\infty$? Calcula-la.

\ex{2} \pau{Cantàbria, extraordinària 2020} Considera $f(x)=\dfrac{1-\cos x}{x}$. Calcula $\displaystyle\lim_{x\to0}f(x)$ i les asímptotes de $f$.

\ex{2} \pau{Canàries, extraordinària 2020} Troba les asímptotes de la funció $f(x)=\dfrac{x^3+5x^2}{x^2-1}$.

\ex{2} \pau{Castella i Lleó, extraordinària 2022} Troba el domini i les asímptotes de $f(x)=\dfrac{x^2}{2-x}$.

\ex{2} \pau{Illes Balears, ordinària 2025} Determina el domini de $f(x)=\dfrac1x\ln(2x)$ i el comportament de la funció en els extrems del domini.

\ex{2} \pau{La Rioja, extraordinària 2023} Troba el domini i les asímptotes de $f(x)=\dfrac{1+4x^4-x^2}{x}$.

\ex{2} \pau{Madrid, extraordinària 2022} Siga $f(x)=\begin{cases}\dfrac{2x+1}{x}&\text{si }x<0\\x^2-4x+3&\text{si }x\geq0\end{cases}$. Estudia la continuïtat de $f$ i calcula, si n'hi ha, les asímptotes horitzontals i verticals.

\ex{2} \pau{País Basc, extraordinària 2025} Troba les asímptotes de $f(x)=\dfrac{x}{x^2-3x-4}$.

\ex{2} \pau{País Basc, ordinària 2024} Troba les asímptotes de $f(x)=\dfrac{x}{x^2-2x+1}$.

\ex{2} \pau{Extremadura, ordinària 2018} Estudia el domini i les asímptotes de $f(x)=\dfrac{1}{x^2-1}$.

\vfill
\begin{center}\color{muted}\small\faCheckCircle[regular]\ \ Ànim! Comprova els teus resultats amb la calculadora de límits dels apunts interactius.\end{center}

\end{document}
